\documentclass[10pt,a4paper]{amsart}
\usepackage[margin=19mm]{geometry}
\usepackage{amsmath,amssymb,mathtools}
\usepackage[scr=rsfs,cal=euler]{mathalfa}
\usepackage{xfrac}
\usepackage[unicode,hidelinks,bookmarks=false]{hyperref}
\newcommand{\ie}{i.e.\ }
\newcommand{\cf}{cf.\ }
\newcommand{\resp}{respectively\ }
\newcommand{\Subsection}{\subsection}
\providecommand{\nobreakdash}{\nobreak\hbox{-}}
\setlength{\emergencystretch}{2em}
\multlinegap0pt
\usepackage{amssymb}
\usepackage{mathtools}
\usepackage[shortlabels]{enumitem}
\usepackage{tikz}
\newtheorem{proposition}{Proposition}
\title{\textbf{On the DML(1) property for regular endomorphisms of affine spaces: the $\mathbb{G}_m$-case}}
\author{She Yang, Aoyang Zheng}
\date{Source: arXiv:2609.29107v1 (CC BY 4.0)}
\begin{document}
\maketitle

\noindent\emph{Source and licence.} \textbf{On the DML(1) property for regular endomorphisms of affine spaces: the $\mathbb{G}_m$-case}. Version arXiv:2609.29107v1 (CC BY 4.0), distributed by arXiv under the Creative Commons Attribution 4.0 International licence, http://creativecommons.org/licenses/by/4.0/, as recorded in the arXiv licence metadata for this exact version and shown on the abstract page https://arxiv.org/abs/2609.29107v1. Full licence notice: \texttt{licenses/arxiv-2609.29107-CC-BY-4.0.txt}.\par\smallskip

\noindent\textit{Editorial scope.} Counted unit: the article's proposition supevolve (source0.tex lines 526--532). The article's complete original proof of it is delivered with it (source0.tex lines 534--582, one \texttt{proof} environment of 49 lines). No mathematics is written, repaired or re-derived: the statement and the proof are the article's own lines, in the article's own notation, and no retained line is substituted or re-flowed.  Retained uncounted as the source-local context the proof needs: lines 126--128; lines 455--463.  Nothing was omitted from any retained span, so the package carries no omission record.  No retained line contains a citation, so the wrapper carries no bibliography and no bibliography entry.  No retained line contains a figure environment, an image-inclusion command or drawing code, so no image is delivered and the package declares no asset dependency.  Editorial formatting only, in addition: an AMS \texttt{amsart} preamble that restates the article's own declarations and macros used by the retained lines, namely the environments \texttt{proposition} on separate counters, together with the standard packages those lines need.  The retained environments print in this document as Proposition 1 in order of appearance, whereas the article prints the same items with the numbering of its own full document.  The complete native source of the article as distributed in the arXiv e-print for this exact version is bundled unchanged as \texttt{sources/211634/source0.tex}.
\begin{proposition}\label{mainprop}
Let $K\subseteq\mathbb{C}$ be an arithmetic function field. Let $f$ be a regular endomorphism of $\mathbb{A}^N$ with $K$-coefficients. Suppose $\deg(f)>1$. Let $(C_{-n})_{n\geq0}$ be a sequence of irreducible closed subcurves of $\mathbb{A}_{K}^N$ which satisfies $f(C_{-n})=C_{-n+1}$ for every $n\geq1$. Suppose the normalization of each $C_{-n}$ is isomorphic to $\mathbb{G}_{m,K}$. Then there are only finitely many different curves in $(C_{-n})_{n\geq0}$. Hence all of them are $f$-periodic.
\end{proposition}

\begin{proposition}\label{bdhtsup}
Let $f$ be the regular endomorphism as in Proposition \ref{mainprop}. For every $n\geq0$, let $P_n\in K[t,t^{-1}]^N$ be the parametrization of $C_{-n}\subseteq\mathbb{A}_K^N$ as above. We have sequences $(a_n)_{n\geq0}\in(K^{\times})^{\mathbb{N}}$ and $(M_n)_{n\geq0}\in\mathbb{Z}_+^{\mathbb{N}}$ such that $f^n(P_n(t))=P_0(a_nt^{M_n})$ holds for every $n$. For each $n$, choose $\lambda_n\in\overline{K}^{\times}$ with $a_n\lambda_n^{M_n}=1$ and put $Q_n(t)=P_n(\lambda_nt)$. Then the following statements hold.
\begin{enumerate}
\item
There exists $H\in\mathbb{R}_{\geq0}$ such that every coefficient of every component of every $Q_n$ has height at most $H$.
\item
There exists $L\in\mathbb{Z}_+$ such that the cardinality of the support of every $P_n$ is at most $L$.
\end{enumerate}
\end{proposition}

\begin{proposition}\label{supevolve}
Let $f$ be the regular endomorphism as above and let $d>1$ be its algebraic degree. Let $L\in\mathbb{Z}_+$ be the number given in Proposition \ref{bdhtsup}(ii). Then there is a finite collection $\mathcal{P}$ of partial linear isometries of $(\mathbb{R}^L,\|\cdot\|_{\infty})$ such that the following holds. For every $n\geq1$, there exists some map $T\colon S\to\mathbb{R}^L$ in $\mathcal{P}$ such that
$$
e(P_n)\in S\quad\text{and}\quad T(e(P_n))=\frac{m_n}{d}e(P_{n-1}).
$$
In particular, we have $m_n\|e(P_{n-1})\|_{\infty}=d\|e(P_n)\|_{\infty}$.
\end{proposition}

\begin{proof}
Let $\mathcal{I}=\{\alpha\in\mathbb{N}^L\mid\|\alpha\|_1\leq d\}$. Let $\mathcal{A}\subseteq M_L(\mathbb{R})=\mathrm{End}(\mathbb{R}^L)$ be the set of matrices that every row is an element of $\mathcal{I}$. Let $\mathcal{S}$ be the set of linear subspaces of $\mathbb{R}^L$ that can be defined by several equations of the form
$$
(\alpha-\beta)\cdot x=0\quad(x\in\mathbb{R}^L)
$$
in which $\alpha,\beta\in\mathcal{I}$. Since $\mathcal{I}$ is finite, so do $\mathcal{A}$ and $\mathcal{S}$. Denote
$$
\mathcal{P}=\left\{\frac{1}{d}A|_{S}\bigg|\ A\in\mathcal{A},\ S\in\mathcal{S},\ \frac{1}{d}A|_{S}\colon S\to\mathbb{R}^L\text{ is an isometry onto its image}\right\}.
$$
We shall show that this finite collection of partial linear isometries satisfies the requirement.

Write $e(P_n)=(e_1,\dots,e_L)$ and write $P_n(t)=\sum\limits_{j=1}^L a_jt^{e_j}$ for some $a_1,\dots,a_L\in K^N$ such that $a_j=0$ for the padded slots $j$. By expanding the expression formally, we find $(B_{\alpha})_{\alpha\in\mathcal{I}}\in(K^N)^{\mathcal{I}}$ such that
$$
f\left(\sum\limits_{j=1}^L a_jt^{k_j}\right)=\sum\limits_{\alpha\in\mathcal{I}}B_{\alpha}t^{\alpha\cdot k}
$$
holds for every $k=(k_1,\dots,k_L)\in\mathbb{Z}^L$.

Partition $\mathcal{I}=I_1\sqcup I_2\sqcup\cdots\sqcup I_t$ according to the value of $\alpha\cdot e(P_n)\ (\alpha\in\mathcal{I})$. Let $S\in\mathcal{S}$ be the linear subspace defined by the equations
$$
(\alpha-\beta)\cdot x=0,\quad\alpha,\beta\in I_j,\quad j=1,\dots,t\quad(x\in\mathbb{R}^L).
$$
Then by definition $e(P_n)\in S$.

We consider a matrix $A\in\mathcal{A}$ defined as follows. Let $l\in\{1,\dots,L\}$. If $l$ is an active slot of $e(P_{n-1})$, then there exists a unique $j_l\in\{1,\dots,t\}$ such that the $l$-th component of $e(P_{n-1})$ is equal to $\frac{\alpha\cdot e(P_n)}{m_n}$ for $\alpha\in I_{j_l}$. We arbitrarily choose an $\alpha\in I_{j_l}$ and let the $l$-th row of $A$ be $\alpha$. If $l$ is a padded slot of $e(P_{n-1})$, then we let the $l$-th row of $A$ be zero.

By definition we have $A\cdot e(P_n)=m_n\cdot e(P_{n-1})$. Therefore, it only remains to check that
$$
T=\frac{1}{d}A|_{S}\colon S\to\mathbb{R}^L
$$
is an isometry onto its image.

Let $\mathcal{H}$ be the finite set of codimension $1$ linear subspaces defined by equations
$$
(\alpha-\beta)\cdot x=0,\quad\alpha,\beta\in\mathcal{I}\text{ lie in different }I_j\quad(x\in\mathbb{R}^L).
$$
Let $U=S\setminus\bigcup\limits_{H\in\mathcal{H}}H$. We have $e(P_n)\in U$, hence $U$ is nonempty. Since the finitely many linear subspaces involved in here are all defined over $\mathbb{Q}$, it suffices to check
$$
\|Tx\|_{\infty}=\|x\|_{\infty}\quad\text{for}\quad x\in U\cap\mathbb{Q}^L.
$$
Pick such an $x=(x_1,\dots,x_L)$ and we may furthermore assume $x\in\mathbb{Z}^L$ by passing to an appropriate multiple. We need to verify $\|Ax\|_{\infty}=d\|x\|_{\infty}$.

By definition, for every $\alpha,\beta\in\mathcal{I}$, we have $(\alpha-\beta)\cdot x=0$ if and only if $(\alpha-\beta)\cdot e(P_n)=0$. Hence for a padded slot $j$ of $e(P_n)$ we have $x_j=0$, and the components of $x$ on the active slots of $e(P_n)$ are pairwise distinct. Write
$$
\tilde{P}(t)=\sum\limits_{j=1}^L a_jt^{x_j}\quad\text{and}\quad Q(t)=f(\tilde{P}(t))=\sum\limits_{\alpha\in\mathcal{I}}B_{\alpha}t^{\alpha\cdot x}.
$$
Then we have $\|e(\tilde{P})\|_{\infty}=\|x\|_{\infty}$. Since $f$ is a regular endomorphism, we also have $\|e(Q)\|_{\infty}=d\|e(\tilde{P})\|_{\infty}$.

When we divide $\mathcal{I}$ according to the value of $\alpha\cdot x\ (\alpha\in\mathcal{I})$, we still get the same partition $\mathcal{I}=I_1\sqcup I_2\sqcup\cdots\sqcup I_t$. This fact, together with the definition of the matrix $A$, lead to the equation $\|e(Q)\|_{\infty}=\|Ax\|_{\infty}$. Hence we get $\|Ax\|_{\infty}=d\|x\|_{\infty}$ and finish the proof.
\end{proof}

\end{document}
